% Chapter 3 -- Working with Windows (Part I, Learning; AN-14 task form).
% Sources (Law 1): os/flair/window.c (document windows, ZoomWindow to the
% standard state (4,40)-(636,476), ConstrainWindowSize 96x64 .. 640x480,
% CollapseWindow); os/milton/kmain.c (flair_live_track_box: a box acts
% only if the button is released inside it; flair_live_do_drag: the grey
% outline and the move on release; a click in an inactive window only
% activates it; flair_live_track_scroll: the arrows repeat while held;
% scroll bars act only in the active window; flair_live_do_finder_close);
% spec/chrome_metrics.h (title bar 22 px, the 12 x 12 boxes, the size box);
% os/flair/finder_windows.h (the root window "Drive A Files" at (20,60)
% 360 x 220, each further window 20 px lower and to the right; at most four
% Finder windows); os/apps/ref_tenant.c (HELLO and NOTES: a click in the
% content shows or hides a marker; the close box ends the application).
% Model: Macintosh User's Guide 1991 Ch. 2 windows sections.
% GATES: test-flair-disk-windows test-flair-drag test-flair-grow test-flair-zoom-toggle test-flair-collapse test-flair-scroll test-flair-close-terminate test-flair-appswitch test-flair-box-track test-window test-pubs-flair-screens
\Gates{test-flair-disk-windows, test-flair-drag, test-flair-grow,
  test-flair-zoom-toggle, test-flair-collapse, test-flair-scroll,
  test-flair-close-terminate, test-flair-appswitch, test-flair-box-track,
  test-window, test-pubs-flair-screens}
\chapter{Working with Windows}

Everything you work with in FLAIR appears in a \term{window}: the
contents of a disk or a folder, a document, a calculator.  This chapter
shows you how to open, arrange and close windows.  The skills are the same
for every window, whichever application it belongs to.

\section{Opening a window}
\index{windows!opening}
\Operation{FINDER-OPEN-VOLUME}
\begin{TaskSteps}
\Step Double-click the \texttt{INITECH} icon.
\Result{A window titled \texttt{Drive A Files} opens.  It shows the
contents of the data disk: the file \file{README.TXT} and the folder
\file{APPS}.}
\end{TaskSteps}

\FlairScreen[label=fig:rootwin]{rootwin}{The Drive A Files window}

The window opens in front of the others, and it is the \term{active
window}: its title bar is striped and its title is black.  The lower menu
bar shows the Finder's menus, because a Finder window is in front.

\section{The parts of a window}
\index{windows!parts of}
Figure~\ref{fig:rootwin} shows the parts every document window has.

\SideHead{Title bar}The bar across the top of the window, with the
window's name in the middle.  Drag the title bar to move the window.

\SideHead{Close box}The small square at the left end of the title bar.
Click it to close the window.\index{close box}

\SideHead{Zoom box}The first of the two boxes at the right end of the title
bar.  Click it to enlarge the window to fill the screen below the menu
bars, and click it again to return the window to its size.\index{zoom box}

\SideHead{Collapse box}The box at the far right of the title bar.  Click it
to roll the window up into its title bar, and click it again to roll the
window down.\index{collapse box}

\SideHead{Scroll bars}The bars along the right side and the bottom of the
window.  When a window is too small to show all of its contents, use the
scroll bars to bring the rest into view.\index{scroll bars}

\SideHead{Size box}The ridged square at the lower right corner.  Drag it to
make the window larger or smaller.\index{size box}

The boxes work like buttons: if you press the mouse button in a box and
change your mind, move the pointer out of the box before releasing the
button, and nothing happens.

\section{Making a window active}
\index{windows!active}
Only one window is active at a time.  To make another window active, click
anywhere in it.
\begin{TaskSteps}
\Step Click the title bar of the \texttt{NOTES} window.
\Result{NOTES comes to the front and becomes the active window.  The window
that was active is drawn inactive: its title bar is plain and its title is
grey.  The lower menu bar changes to NOTES's menus.}
\end{TaskSteps}

\FlairScreen[label=fig:activate]{activate}{NOTES made active}

A click in an inactive window only brings it to the front.  It does not
also act on what is under the pointer: click again to do that.

\section{Moving a window}
\index{windows!moving}
\Operation{FLAIR-DRAG}
\begin{TaskSteps}
\Step Point at the title bar of the \texttt{Drive A Files} window, away
from the boxes, and drag.
\Result{An outline of the window follows the pointer.}
\Step Release the button.
\Result{The window moves to where the outline is.  FLAIR remembers where
you put each Finder window and opens it there next time.}
\end{TaskSteps}

\FlairScreen[label=fig:moved]{moved}{The Drive A Files window, moved}

\section{Changing the size of a window}
\index{windows!changing size}
\Operation{FLAIR-GROW}
\begin{TaskSteps}
\Step Point at the size box at the lower right corner of the window and
drag it down and to the right.
\Result{An outline shows the new size of the window.}
\Step Release the button.
\Result{The window takes the size of the outline.}
\end{TaskSteps}

\FlairScreen[label=fig:grown]{grown}{The Drive A Files window, enlarged}

A window cannot be made smaller than 96 dots wide and 64 high, and cannot
be made larger than the screen.

\section{Zooming a window}
\index{zooming}
\Operation{FLAIR-ZOOM}
\begin{TaskSteps}
\Step Click the zoom box of the \texttt{Drive A Files} window.
\Result{The window fills the screen below the menu bars.}
\end{TaskSteps}

\FlairScreen[label=fig:zoomed]{zoomed}{The Drive A Files window, zoomed}

\begin{TaskSteps}[resume]
\Step Click the zoom box again.
\Result{The window returns to the size and place it had.  If you change
the size of a zoomed window with its size box, the new size becomes the
window's own size.}
\end{TaskSteps}

\section{Collapsing a window}
\index{collapsing a window}
\Operation{FLAIR-COLLAPSE}
\begin{TaskSteps}
\Step Click the collapse box of the \texttt{Drive A Files} window.
\Result{The window rolls up until only its title bar is left.  You can
still move it by its title bar.}
\end{TaskSteps}

\FlairScreen[label=fig:collapsed]{collapsed}{The Drive A Files window,
collapsed}

\begin{TaskSteps}[resume]
\Step Click the collapse box again.
\Result{The window rolls down to its full size.}
\end{TaskSteps}

\section{Scrolling}
\index{scrolling}
\Operation{FLAIR-SCROLL}
A window often holds more than it can show.  The \file{APPS} folder holds
twenty files, and its window shows only the first rows of them.
\begin{TaskSteps}
\Step Double-click the \file{APPS} folder in the \texttt{Drive A Files}
window.
\Result{The \texttt{APPS} window opens in front.}
\Step Click the arrow at the bottom of the scroll bar on the right side of
the window.
\Result{The contents move up so that more of them come into view.  The
small square in the scroll bar, the \term{scroll box}, moves down to show
where you are.}
\end{TaskSteps}

\FlairScreen[label=fig:scrolled]{scrolled}{The APPS window, scrolled}

There are three ways to scroll:
\begin{itemize}
\item Click an arrow to move a little.  Hold the button down to keep
  moving.
\item Click the grey part of the scroll bar, above or below the scroll
  box, to move by a windowful.
\item Drag the scroll box to move to a place in proportion: to the end of
  the contents, drag it to the end of the bar.
\end{itemize}
Scroll bars work only in the active window; a click in the scroll bar of an
inactive window makes the window active.

\section{Closing a window}
\index{windows!closing}
\Operation{FINDER-CLOSE-WINDOW}
\begin{TaskSteps}
\Step Click the close box of the \texttt{APPS} window.
\Result{The window closes.  The next window behind it becomes the active
window.}
\end{TaskSteps}

\FlairScreen[label=fig:closed]{closed}{The APPS window closed; Drive A Files
is active again}
When the Finder is the active application you can also close the front
Finder window by choosing \menuitem{Close Window} from the \menuname{File}
menu, or by pressing \key{Ctrl-W}.

Closing a Finder window does not change anything on the disk.  Closing a
window of an application usually ends the application; an application
with unsaved work asks you first (Chapter~5).

\section{Practising with HELLO and NOTES}
\index{HELLO}\index{NOTES}
The HELLO and NOTES windows that FLAIR opens when it starts are there for
practice.  Move them, change their size, zoom and collapse them, and bring
each to the front in turn by clicking it.  Each has its own menu bar, which
appears in the lower bar while its window is in front (Chapter~8 lists
their menus).

When you close HELLO or NOTES with its close box, the window and its
application are gone until you start FLAIR again.

\FlairScreen[label=fig:helloclosed]{helloclosed}{HELLO closed}
\Operation{FLAIR-CLOSE}
